% =====================================================================
%  COMP3320 2026 -- Lab 3 interactive guide slides
%  Source of truth: script.md  (this file is the LaTeX translation)
%  Theme: cookie (LuaLaTeX-first).  Build with:  latexmk slides.tex
%  (or:  lualatex slides.tex  twice, if latexmk is unavailable)
%
%  The deck is the session companion for the lab README: every beat is
%  predict -> run -> reveal, so most frames carry \pause overlays.
%  Nothing here spoils a reference number; the reveals are the ones the
%  README itself states in prose.
% =====================================================================
\documentclass[aspectratio=169]{beamer}

\usepackage{amsmath,amssymb}
\usepackage{booktabs}
\usepackage{tabularx}

% Logos etc. live in the shared template assets; add local assets/ too.
\graphicspath{{assets/}{../../template/assets/}}

\definecolor{fortranPurple}{HTML}{503874}     % course deep purple
\definecolor{fortranPurpleMid}{HTML}{7E5AA6}
\usetheme[
  accent=fortranPurple,
  progressbar=foot,
  sectionpage=progressbar,
  subsectionpage=none,
  numbering=fraction,
  numberrail=section,
  block=fill,
  notes=hide,
  toc=aftertitle,
  closing=contact
]{cookie}

% small muted right-aligned annotation, as in the lab 1 deck
\newcommand{\cplx}[1]{{\footnotesize\color{cookieMuted}#1}}

\title[COMP3320 Lab 3]{Programming with SIMD}
\subtitle{What the vector unit can do that the compiler will not}
\author[COMP3320 2026 \textperiodcentered\ Lab 3]{Jorge Luis G\'alvez Vallejo}
\cookieemail{jorge.galvezvallejo@anu.edu.au}
\institute{School of Computing}
\cookieuni{Australian National University}
\cookielocation{Canberra, Australia}
\date{Week 4, 2026}
\cookietagline{COMP3320 \textperiodcentered\ Lab 3 \textperiodcentered\ Week 4}

\cookieclosingtitle{}
\cookieclosingtext{Predict. Measure. Argue with the result.}

\begin{document}

\maketitle

% =====================================================================
\section{The story you have been told}

\begin{frame}{Every optimisation guide says the same thing}
  \begin{cookiequote}
    The compiler cannot vectorise \texttt{a[i] = a[i] + b[i]}, because
    \texttt{a} and \texttt{b} might overlap. Add \texttt{restrict} to promise
    that they do not, and it will vectorise.
  \end{cookiequote}
  \bigskip
  \pause
  This lab \textbf{tests} that claim instead of believing it.
  \begin{itemize}\setlength\itemsep{4pt}
    \item<3-> It is \textbf{largely false} on a modern compiler.
    \item<4-> The loop that genuinely refuses does so for a reason that is
      \textbf{invisible in the source}.
    \item<5-> The vector unit \emph{can} do conditional work --- the hardware has
      existed since 2016 --- but \textbf{C has no way to say so}.
      So the compiler cannot use it, and you have to.
  \end{itemize}
\end{frame}

\begin{frame}{By hand no longer means intrinsics}
  This lab uses the \textbf{Vector Class Library} (VCL), vendored in the repo ---
  nothing to download, nothing to \texttt{module load}.
  \medskip
  \begin{itemize}\setlength\itemsep{4pt}
    \item Same instructions as intrinsics --- you will \emph{check} that in Q11.
    \item Readable: \texttt{a + b} and \texttt{select(mask, x, y)},
      not \texttt{\_mm512\_mask\_mov\_pd}.
    \item Portable: the same source compiles for SSE2, AVX2 or AVX-512,
      and \emph{degrades} to narrower vectors instead of failing to build.
  \end{itemize}
  \bigskip
  \begin{block}{The problem, all lab long}
    A \textbf{dam break}: a wall is removed, a wave runs over dry ground.
    The wet/dry front \emph{moves} --- so the inner loop has a branch whose
    outcome genuinely changes as the simulation runs.
  \end{block}
\end{frame}

\begin{frame}{How to use the lab}
  \textbf{Longer than the session, and none of it is marked.} Both deliberate.
  \medskip
  \begin{columns}[T,onlytextwidth]
    \column{.48\textwidth}
      \begin{cookiecard}[title={In the session (core path)}]
        Questions \textbf{1, 2, 3, 5, 7, 9, 10, 17}.\\[2pt]
        \cplx{Predict, catch the folklore, measure the three routes, choose
        one, look back.}
      \end{cookiecard}
    \column{.48\textwidth}
      \begin{cookiecard}[title={In your own time}]
        \textbf{Q13 first} --- the harness checks your kernel bit for bit.\\
        Then \textbf{Q11}, the deepest question.\\[2pt]
        \cplx{\emph{Extension.} = drop if short of time. Part 5 is optional
        --- and the closest thing here to the actual job.}
      \end{cookiecard}
  \end{columns}
  \bigskip
  \pause
  Several of the answers ahead are \textbf{not what you expect} --- that is why
  this is worth doing with a demonstrator in the room to argue with.
\end{frame}

% =====================================================================
\section{Setup: three traps before you start}

\begin{frame}[fragile]{Trap 1: get a compute node, and do not waste it}
  A login node once gave \texttt{9.28 3.30 8.57 9.28 8.56} ns/element for the
  \emph{same command}.
  \begin{cookiecode}[title={ask for this}]{bash}
qsub -I -q normal -l ncpus=1,mem=4GB,walltime=2:00:00,storage=scratch/$PROJECT -P $PROJECT
  \end{cookiecode}
  \pause
  \begin{itemize}\setlength\itemsep{3pt}
    \item \textbf{Why 4\,GB?} Gadi charges the larger of CPUs and memory, one
      charged CPU per 4\,GB --- \texttt{mem=16GB} costs \textbf{4$\times$} for
      identical work. The whole lab is ${\sim}2$ SU/hour.
    \item \textbf{The wrong node fails \emph{silently}.} VCL emulates
      \texttt{Vec8d} on Broadwell: everything builds and runs --- and none of
      your numbers mean anything.
  \end{itemize}
  \begin{cookiecode}[title={check when you land (and Q1 makes you print INSTRSET)}]{bash}
lscpu | grep -oE 'avx512[a-z0-9_]*' | sort -u
  \end{cookiecode}
\end{frame}

\begin{frame}[fragile]{Trap 2: the environment}
  \begin{cookiecode}[title={not optional --- the lab will not build without this}]{bash}
module load gcc/15.1.0
module unload nvidia-hpc-sdk
unset CPATH C_INCLUDE_PATH CPLUS_INCLUDE_PATH
  \end{cookiecode}
  \medskip
  \begin{itemize}\setlength\itemsep{4pt}
    \item nvidia-hpc-sdk ships its own \texttt{x86intrin.h}; VCL includes it;
      the build dies in \textbf{hundreds of lines from inside \texttt{/apps/}}
      that never mention your code. \texttt{make} checks and refuses to start.
    \item The system gcc is from \textbf{2018}. Every reference number assumes
      15.1.0, and several answers genuinely differ between the two.
      \\ \cplx{Questions 5 and 6 go back to the old compiler on purpose --- and
      say so.}
  \end{itemize}
\end{frame}

\begin{frame}[fragile]{Trap 3: pin every timing run}
  \begin{cookiecode}[title={define pin() once --- every timing command uses it}]{bash}
pin() {
    local cpu=$(awk '/Cpus_allowed_list/{split($2,a,",");
                     split(a[1],b,"-");print b[1]}' /proc/self/status)
    numactl --localalloc --physcpubind=$cpu "$@"
}
  \end{cookiecode}
  \cplx{Binds to the first CPU your job actually owns, memory on that socket.}
  \medskip
  \pause
  \begin{itemize}\setlength\itemsep{4pt}
    \item Unpinned, these measurements varied by \textbf{3$\times$} --- not
      noisy, \emph{wrong}, and wrong in the wrong direction.
    \item Even pinned: run everything more than once and \textbf{know your own
      scatter}. Two numbers inside it are \emph{the same number} --- and ``I
      could not resolve this'' is a correct, complete answer.
  \end{itemize}
\end{frame}

% =====================================================================
\section{Part 1: the vocabulary}

\begin{frame}[fragile]{Eight velocities at once --- but only in the wet cells}
  \begin{cookiecode}[title={intrinsics: correct, fast, unreadable, AVX-512 only}]{C}
__mmask8 wet = _mm512_cmp_pd_mask(hL, dry, _CMP_GT_OQ);
__m512d  uL  = _mm512_maskz_div_pd(wet, huL, hL);
  \end{cookiecode}
  \pause
  \begin{cookiecode}[title={the same two lines in VCL}]{C++}
Vec8db wet = hL > H_DRY;
Vec8d  uL  = select(wet, huL / hL, 0.0);
  \end{cookiecode}
  \pause
  \begin{block}{A boolean vector is the whole idea}
    A comparison on \texttt{Vec8d} gives \textbf{eight bools} in a hardware mask
    register; \texttt{select} takes each lane from one side or the other.
    That pair is how a vector unit does conditional work --- and C cannot
    write it. That sentence is the reason this lab exists.
  \end{block}
\end{frame}

\begin{frame}{Two operations to respect now, because they bite later}
  \begin{itemize}\setlength\itemsep{10pt}
    \item \texttt{load\_partial(m, p)} --- loads $m$ elements,
      \textbf{zeroes the rest}. The library's answer to a loop tail.
      \\ \cplx{Only harmless if zero is a harmless input to what you compute.
      Q13c asks exactly when the zero fill \emph{lies}.}
    \item \texttt{horizontal\_max(a)} --- eight lanes $\to$ one scalar.
      The only listed operation that is \textbf{not} one instruction per lane:
      a sequence of shuffles and maxes.
      \\ \cplx{Where you put it decides whether your loop is fast.
      That is the whole of Q13.}
  \end{itemize}
  \bigskip
  Everything else you need --- loads, stores, arithmetic, \texttt{select},
  \texttt{if\_add}, horizontal reductions --- is one table in the README.
\end{frame}

\begin{frame}[standout,noframenumbering]
  Question 2: read the four variants,\\ then \textbf{write your prediction down}.\\[8pt]
  {\normalsize How much faster will the hand-written variant 2 be, and why?\\
  Question 17 collects this debt --- show a demonstrator now.}
\end{frame}

% =====================================================================
\section{Part 2: the folklore}

\begin{frame}[fragile]{Poll: will gcc \texttt{-O3} vectorise this?}
  \begin{cookiecode}[title={compute.c --- no restrict anywhere}]{C}
void compute(int n, double *a, double *b)
{
  for (int i = 0; i < n; i ++) {
    a[i] = a[i] + b[i];
  }
}
  \end{cookiecode}
  \medskip
  Hands up: \textbf{vectorised}, or \textbf{refused because \texttt{a} and
  \texttt{b} might overlap}?
  \pause
  \begin{block}{Don't take a vote's word for it either --- ask the compiler (Q3)}
    \texttt{-fopt-info-vec} on both files, then sweep
    \{gcc 8.5, gcc 15.1\} $\times$ \{\texttt{-O2}, \texttt{-O3}\}.
    In \textbf{three} of the four configurations \texttt{restrict} changes
    nothing at all. In \textbf{one} it is decisive --- and not the one you
    would guess.
  \end{block}
\end{frame}

\begin{frame}[fragile]{So what does \texttt{restrict} actually decide? (Q5)}
  The driver calls \texttt{compute(nsize-1, \&a[1], \&a[0])} --- the same array,
  offset by one. With \texttt{restrict}, that promise is now \textbf{a lie}.
  \begin{cookiecode}[title={watch it go wrong}]{bash}
module unload gcc; module load gcc/system
make -B exe_runcompute exe_runcompute_restrict
./exe_runcompute_restrict 8      # then try 12
  \end{cookiecode}
  \pause
  The output is right for small inputs, then wrong --- \textbf{in repeating
  groups}. Poll: what sets the group size?
  \pause
  \begin{itemize}\setlength\itemsep{3pt}
    \item Not \texttt{restrict} itself: the \textbf{width of the instructions}
      the compiler emitted. \texttt{objdump} shows you which register class.
    \item No warning, no crash --- a \emph{silently wrong answer}.
      So: when \emph{should} you use \texttt{restrict}?
  \end{itemize}
\end{frame}

\begin{frame}[fragile]{The loop that refuses anyway (Q6)}
  \texttt{vecblock.c}: every pointer \texttt{restrict}, three separate
  allocations, \textbf{no aliasing at all} --- and the report still says
  \begin{center}
    \texttt{not vectorized: control flow in loop.}
  \end{center}
  \pause
  Poll: variant 1 contains \textbf{no \texttt{if} statement}. Where is the
  control flow?
  \pause
  \begin{block}{The branch you did not write}
    \texttt{sqrt} is obliged to set \texttt{errno} for a negative argument ---
    a hidden branch in a maths call. \texttt{-fno-math-errno} removes the
    obligation, and the loop vectorises.
  \end{block}
  \cplx{Q6e: gcc 15 proves the argument non-negative by itself and vectorises
  anyway --- everything above is a statement about \emph{one compiler, from
  2018}. What is the shelf life of ``the compiler won't do X''?}
\end{frame}

% =====================================================================
\section{Part 3: the measurement}

\begin{frame}[fragile]{Four variants, one answer (Q7)}
  \begin{cookiecode}[title={back on gcc/15.1.0 first --- then fill the table}]{bash}
module unload gcc; module load gcc/15.1.0
make -B exe_swe exe_swe_fast
pin ./exe_swe 1000000 <variant> 100
  \end{cookiecode}
  \begin{center}\small
    \begin{tabular}{@{}lcccc@{}}
      \toprule
      ns per face & 0 scalar \texttt{if} & 1 branchless & 2 VCL & 3 restructured \\
      \midrule
      \texttt{exe\_swe}       & ? & ? & ? & ? \\
      \texttt{exe\_swe\_fast} & ? & ? & ? & ? \\
      total mass              & \multicolumn{4}{c}{\emph{every digit, both builds}} \\
      \bottomrule
    \end{tabular}
  \end{center}
  \medskip
  Mass is conserved \emph{exactly} by the scheme --- the printed total is a real
  correctness check, not a formality.
  \pause
  \cplx{The \texttt{-ffast-math} build's mass is \emph{closer} to the exact
  500000. So what, precisely, is still the problem with that flag?}
\end{frame}

\begin{frame}[fragile]{Poll: variants 0 and 1 run in \emph{identical} time. Why? (Q8)}
  Not similar. Identical, run after run. That is not noise.
  \pause
  \begin{cookiecode}[title={ask the binary}]{bash}
nm exe_swe | grep swe_flux
objdump -d exe_swe | grep -A2 '<swe_flux_branchless>:'
  \end{cookiecode}
  \pause
  \begin{block}{The compiler folded them into one function}
    It proved your ``clever branchless rewrite'' equivalent to the plain
    \texttt{if} and kept one copy. Rewriting a branch by hand and timing both
    versions in the same binary can be measuring \textbf{nothing at all}.
  \end{block}
\end{frame}

\begin{frame}{Three routes to a vectorised flux loop (Q9, Q10)}
  The refusal is the same string as Q6 --- this time the hidden control flow is
  a \textbf{conditional divide} (what may a divide do that a multiply cannot?).
  \medskip
  \begin{columns}[T,onlytextwidth]
    \column{.32\textwidth}
      \begin{cookiecard}[title={write it yourself}]
        VCL, variant 2. Works. Costs: code you must own, a dependency, and
        \emph{knowing it matters}.
      \end{cookiecard}
    \column{.32\textwidth}
      \begin{cookiecard}[title={restructure source}]
        Variant 3: select the \emph{operands}, divide unconditionally.
        Vectorises --- measure what it buys before you admire it.
      \end{cookiecard}
    \column{.32\textwidth}
      \begin{cookiecard}[title={\texttt{-ffast-math}}]
        One flag, whole program. You did not ask for the arithmetic to change,
        and cannot tell where it did.
      \end{cookiecard}
  \end{columns}
  \bigskip
  \pause
  \textbf{Which do you ship} in a code whose defining requirement is that the
  answer does not change?
\end{frame}

\begin{frame}[fragile]{Q11: eight lanes, nowhere near $8\times$. Why?}
  First prove VCL gave you the instructions you asked for
  (\texttt{make asm}: count \texttt{vdivpd}/\texttt{vsqrtpd}, find no
  \texttt{call}). Then: they are \textbf{not fully pipelined} --- estimate the
  loop's floor from their throughput, and build the roofline (DRAM ceiling from
  \texttt{exe\_alias}, peak from FMA units $\times$ width $\times$ clock).
  \pause
  \begin{cookiecode}[title={the size sweep --- note what nsteps is doing}]{bash}
for n in 2000 20000 200000 2000000; do
    pin ./exe_swe $n 2 $((200000000/n)) | grep ns/face
done
  \end{cookiecode}
  \begin{block}{Hold total work fixed when you sweep}
    Fix \texttt{nsteps} instead and the small sizes finish before the clock
    leaves idle --- a ramp-up that looks \emph{exactly} like a cache curve.
  \end{block}
\end{frame}

% =====================================================================
\section{Afterwards, and the real world}

\begin{frame}{Q13: write one yourself --- start here after the session}
  The CFL reduction $s_{\max} = \max_i\,(|u_i| + \sqrt{g\,h_i}\,)$:
  $n$ values in, \textbf{one} out. Structurally not the flux loop --- every lane
  must reach one scalar, and \emph{how} is the whole exercise.
  \medskip
  \begin{itemize}\setlength\itemsep{4pt}
    \item The harness checks your kernel \textbf{bit for bit} --- including
      $n \not\equiv 0 \pmod 8$ and dry cells. You need nobody to mark it.
    \item The three classic mistakes are listed in \texttt{maxwave.cpp}.
      \cplx{Where the horizontal collapse goes; what the tail zeros poison;
      what a dry cell divides by.}
    \item Doing it yourself is the point --- the README asks AI assistants not
      to write this one for you, and the harness knows if they did not listen.
  \end{itemize}
\end{frame}

\begin{frame}{Part 5: the production kernel (optional, and the actual job)}
  A real HLL Riemann solver from an ocean model, in \textbf{Fortran} --- the
  same three obstacles.
  \medskip
  \begin{itemize}\setlength\itemsep{4pt}
    \item \textbf{Q15 first if time is short}: hand-inlining the per-cell helper
      is everyone's first hypothesis. It changes exactly one of \{time, report\}
      --- find out which, and what that separates.
    \item \textbf{Q16 is the deliverable}: one day, answer must not move beyond
      rounding, other compilers must keep working. Write the case for and
      against rewriting.
  \end{itemize}
  \medskip
  \begin{cookiequote}
    ``Here is the bottleneck, here is why, here is the fix, here is what it is
    worth --- and here is why I am not doing it.''
  \end{cookiequote}
\end{frame}

\begin{frame}[standout,noframenumbering]
  Question 17:\\ take out your question 2 prediction.\\[8pt]
  {\normalsize How wrong were you --- and \emph{where}?}
\end{frame}

% closing slide is inserted automatically by `closing=contact`
\end{document}
